\documentclass[aspectratio=169,10pt,fontset=fandol]{ctexbeamer}
\usepackage{amsmath,amssymb,mathtools,bm,mathrsfs}
\usepackage{graphicx,booktabs,array,tabularx,multirow}
\usepackage{tikz,pgfplots,xparse,xstring,mdframed}
\usetikzlibrary{arrows.meta,calc,angles,positioning,shapes.geometric,graphs}
\pgfplotsset{compat=1.18}
\usetheme{fduclassroom}
\setbeamerfont{normal text}{size=\normalsize}
\setbeamerfont{title}{size=\large}
\setbeamerfont{subtitle}{size=\normalsize}
\setbeamerfont{frametitle}{size=\large}
\setbeamerfont{framesubtitle}{size=\small}
\setbeamerfont{block title}{size=\normalsize}
\setbeamerfont{block body}{size=\normalsize}
\setbeamerfont{section in toc}{size=\normalsize}
\setbeamerfont{section in head/foot}{size=\footnotesize}
\setbeamerfont{author in head/foot}{size=\footnotesize}
\setbeamerfont{title in head/foot}{size=\footnotesize}
\setbeamerfont{date in head/foot}{size=\footnotesize}
\setbeamertemplate{navigation symbols}{}
\setbeamertemplate{itemize item}{\normalsize\textbullet}
\setbeamertemplate{itemize subitem}{\small\textbullet}
\setbeamercolor{structure}{fg=blue!65!black}
\setbeamercolor{block title}{fg=blue!65!black,bg=blue!8}
\setbeamercolor{block body}{fg=black,bg=blue!3}

% ctex fontset=fandol; system fonts are not redistributed.


\title{数列极限}
\begin{document}
\section{数列极限}
% FRAME-ID F001
\begin{frame}[t,fragile]{目录}
\normalsize
{\color{blue}\normalsize 1. 数列极限\par}\vspace{0.3em}
{\color{gray!40}\normalsize 2. 极限唯一性\par}\vspace{0.3em}
\end{frame}
% FRAME-ID F002
\begin{frame}[t,fragile]{数列极限}
\normalsize
\begin{block}{定义}
{\normalsize
% SOURCE-BEGIN U0001 0 86
若对每个 $\varepsilon>0$，存在 $N\in\mathbb{N}$，当 $n>N$ 时有 $|a_n-a|<\varepsilon$，则称数列收敛于 $a$。
% SOURCE-END U0001
}
\end{block}
\end{frame}
% FRAME-ID F003
\begin{frame}[t,fragile]{数列极限}
\normalsize
\begin{block}{例题}
{\normalsize
% SOURCE-BEGIN U0002 0 30
计算 $\lim_{n\to\infty}\frac1n$。
% SOURCE-END U0002
}
\end{block}
\pause
{\normalsize
\textbf{解：}
% SOURCE-BEGIN U0003 0 79
任取 $\varepsilon>0$，取 $N>1/\varepsilon$，当 $n>N$ 时 $|1/n|<\varepsilon$，所以极限为 $0$。
% SOURCE-END U0003
}
\end{frame}
\section{极限唯一性}
% FRAME-ID F004
\begin{frame}[t,fragile]{目录}
\normalsize
{\color{black}\normalsize 1. 数列极限\par}\vspace{0.3em}
{\color{blue}\normalsize 2. 极限唯一性\par}\vspace{0.3em}
\end{frame}
% FRAME-ID F005
\begin{frame}[t,fragile]{极限唯一性}
\normalsize
\begin{block}{定理}
{\normalsize
% SOURCE-BEGIN U0004 0 10
收敛数列的极限唯一。
% SOURCE-END U0004
}
\end{block}
\pause
{\normalsize
\textbf{证明：}
% SOURCE-BEGIN U0005 0 68
若两个极限为 $a\ne b$，取 $\varepsilon=|a-b|/3$，由三角不等式得 $|a-b|<2|a-b|/3$，矛盾。
% SOURCE-END U0005
}
\end{frame}
\end{document}
